\uuid{3L7D}
\titre{Famille libre/Génératrice/Polynômes}
\auteur{hanani ; mzari}
\organisation{exo7}
\licence{CC BY-NC-SA 4.0}
\datecreate{2020-01-20}
\niveau{L1}
\module{Algèbre}
\chapitre{Espace vectoriel}
\sousChapitre{Base}
\difficulte{2}
\theme{Base et dimension, Famille libre, Génératrice, Polynômes}
\exo7qcmid{lille-2.espaces-vectoriels.031}

\contenu{
\begin{question}
\qclassification{106.04}
\qtags{motcle=Famille libre/Génératrice/Polynômes}
Dans $\Rr_1[X]$, l'ensemble des polynômes à coefficients réels de degré $\le 1$, on considère les polynômes $P_1= X+1, P_2= X-1, P_3 = 1$. Quelles sont les assertions vraies ?
\begin{answers}
\bad{$\{P_1,P_2,P_3\}$ est une famille libre}.
\good{$\{P_1,P_2,P_3\}$ est une famille génératrice de $\Rr_1[X]$}.
\bad{$\{P_1,P_2,P_3\}$ est une base de $\Rr_1[X]$}.
\good{$\{P_2,P_3\}$ est une base de $\Rr_1[X]$}.
\end{answers}
\begin{explanations} On a : $P_1-P_2-2P_3=0$, donc $\{P_1,P_2,P_3 \}$ n'est pas libre. Par contre, $\{P_1,P_2,P_3 \}$ est une famille génératrice de $\Rr_1[X]$ puisqu'elle contient $2$ polynômes non colinéaires. Toute famille extraite de $\{P_1,P_2,P_3 \}$, contenant $2$ vecteurs, est une base de $\Rr_1[X]$.
\end{explanations}
\end{question}
}
